\section*{अध्याय \ref{ch:b3:electrons-in-solids} --- ठोसों में इलेक्ट्रॉन}

\begin{solution}{exo:b3:electrons-in-solids:1}
(क) $\tau = m_{\text{e}}/ne^2\rho = \qty{2.5e-14}{s}$। (ख) $\ell =
v_{\text{F}}\tau \approx \qty{39}{nm}$। (ग) लगभग 110 \omterm{def:b3:crystalline-solids:lattice}{जालक
नियतांक}। (घ) सौ आयनों के पास से बिना प्रकीर्ण हुए गुज़रता कोई
इलेक्ट्रॉन उन्हीं आयनों से टकरा नहीं सकता ---
\omterm{def:b3:crystalline-solids:lattice}{जालक} को उसके लिए पारदर्शी होना ही चाहिए, और ठीक यही \omterm{thm:b3:electrons-in-solids:bloch}{ब्लॉख की प्रमेय}
(\cref{thm:b3:electrons-in-solids:bloch}) आगे चलकर सिद्ध करती है।
\end{solution}

\begin{solution}{exo:b3:electrons-in-solids:2}
(क) $v_{\text{d}} = I/neA = \qty{7.3e-4}{m/s}$ --- यानी प्रति सेकंड
एक मिलीमीटर से भी कम। (ख) लगभग 23 मिनट। (ग) वह \emph{क्षेत्र}
तार के अनुदिश लगभग प्रकाश की गति से जम जाता है और पूरे इलेक्ट्रॉन
सागर को एक साथ बहने लगाता है: चलने वाले धीमे हैं,
चलने का आदेश नहीं। (घ) $v_{\text{d}}/v_{\text{F}} \sim
10^{-9}$: चालन किसी प्रचंड क्वांटम गति पर लगा एक अगोचर झुकाव है, न
कि कोई शांत चिरसम्मत बहाव।
\end{solution}

\begin{solution}{exo:b3:electrons-in-solids:3}
(क) आधी भरी बैंड: \omterm{def:b3:electrons-in-solids:classes}{धातु}, $\dd\rho/\dd T > 0$। (ख) दो
इलेक्ट्रॉन उसकी बैंड ठीक-ठीक भर देते; मैग्नीशियम इसलिए चालन करता है
कि उसकी भरी बैंड अगली खाली बैंड से \emph{अतिव्यापी} है: \omterm{def:b3:electrons-in-solids:classes}{धातु},
$\dd\rho/\dd T > 0$। (ग) भरी बैंड, \qty{5.5}{eV} अंतराल: \omterm{def:b3:electrons-in-solids:classes}{विद्युत्रोधी}
(औपचारिक रूप से $\dd\rho/\dd T < 0$, पर गिनने को वाहक हैं ही
नहीं)। (घ) \omterm{def:b3:electrons-in-solids:classes}{अर्धचालक}: $\dd\rho/\dd T < 0$, यानी वही
चरघातांकी चिह्न-पलटाव, जो उसके वर्ग की चुगली कर देता है।
\end{solution}

\begin{solution}{exo:b3:electrons-in-solids:4}
(क) $\lambda = hc/E_{\text{g}} = \qty{1.1}{\micro m}$: उससे आगे
सिलिकॉन अवरक्त में पारदर्शी है --- और इसीलिए अवरक्त
कैमरों के लेंस सिलिकॉन के बन सकते हैं। (ख) \qty{225}{nm}, यानी
पराबैंगनी में: कोई दृश्य फ़ोटॉन अवशोषित हो ही नहीं सकता, इसलिए हीरा
पानी-सा साफ़ है। (ग) \qty{365}{nm}: नीले (\qty{2.8}{eV}) के लिए
पर्याप्त चौड़ा अंतराल नाइट्राइडों में है और व्यवहार में लगभग कहीं और
नहीं --- और इसीलिए नीले LED ने पदार्थ-वृद्धि पर दशकों प्रतीक्षा की।
(घ) उसका अंतराल ही उसका फ़ोटॉन तय करता है: अधिक चलाने से ऊष्मा बढ़ती
है, अंतराल की चौड़ाई नहीं, और वह डायोड अँधेरे में ही पक जाता है ---
किसी LED का रंग पदार्थ का नियतांक है, कोई चमक-सेटिंग नहीं।
\end{solution}

\begin{solution}{exo:b3:electrons-in-solids:5}
(क) $\kappa/\sigma T = 400/(5.9\times10^7\times300) =
\qty{2.3e-8}{W.\ohm/K^2}$। (ख) $L$ के दस प्रतिशत भीतर। (ग)
वही फ़र्मी-सतह वाले इलेक्ट्रॉन आवेश और ऊष्मा, दोनों ढोते हैं, अतः
उनका प्रकीर्णन-काल उस अनुपात में कट जाता है। (घ) इस्पात का बर्तन
भोजन तक ऊष्मा इसलिए ले जाता है कि उसके इलेक्ट्रॉन चलते हैं;
और लकड़ी का हत्था, जो दोनों अर्थों में \omterm{def:b3:electrons-in-solids:classes}{विद्युत्रोधी} है, उन्हें ---
और ऊष्मा को --- आपके हाथ से दूर रखता है।
\end{solution}

\begin{solution}{exo:b3:electrons-in-solids:6}
(क) पचास प्रतिशत अतिरिक्त --- जो प्रयोग में खुल्लम-खुल्ला अनुपस्थित है।
(ख) $T/T_{\text{F}} \approx 0.004$: इलेक्ट्रॉनी पद
एक प्रतिशत से भी नीचे गला दिया जाता है। (ग) \omterm{def:b3:crystalline-solids:lattice}{जालक} का $T^3$
इलेक्ट्रॉनों के $T$ से तेज़ी से ढहता है: कुछ ही केल्विन से नीचे वही
``नगण्य'' इलेक्ट्रॉन बाकी बचते हैं। (घ) $C/T =
\gamma + \beta T^2$: अंतःखंड $\gamma$ इलेक्ट्रॉनों को तौलता है,
और ढाल $\beta$ \omterm{def:b3:crystalline-solids:lattice}{जालक} को --- एक ही आलेख, दोनों किरायेदार।
\end{solution}

\begin{solution}{exo:b3:electrons-in-solids:7}
(क) $k_{\text{F}} = (3\pi^2n)^{1/3} = \qty{1.36e10}{m^{-1}}$:
$\lambda_{\text{F}} = \qty{4.6}{\angstrom}$। (ख) $2a =
\qty{7.2}{\angstrom}$: वही कोटि --- ताँबे की फ़र्मी सतह
क्षेत्र-सीमा की ओर पहुँच जाती है, और अंतराल खोलने वाली भौतिकी ठीक
\emph{उन्हीं} ऊर्जाओं पर होती है जो मायने रखती हैं। (ग) एक अप्रगामी
तरंग अपना आवेश धन आयनों पर ढेर करती है (कम स्थिरवैद्युत
ऊर्जा), और दूसरी उनके बीच (अधिक): वही तरंगदैर्घ्य, दो
ऊर्जाएँ --- यानी वह अंतराल। (घ) 110 जालक-नियतांक वाला मुक्त पथ:
ब्लॉख तरंगें किसी निर्दोष \omterm{def:b3:crystalline-solids:lattice}{जालक} से प्रकीर्ण नहीं होतीं, अतः प्रतिरोध
के लिए केवल कंपन और अशुद्धियाँ बचती हैं।
\end{solution}

\begin{solution}{exo:b3:electrons-in-solids:8}
(क) $E_{\text{g}}/2k_{\text{B}} = \qty{6500}{K}$:
$\exp[6500(1/300 - 1/400)] \approx 230$। (ख) थर्मिस्टर का
प्रतिरोध चरघातांकी ढंग से गोता लगाता है --- एक तीखा, संवेदी वक्र;
जबकि प्लैटिनम का धीरे-धीरे और रैखिक रूप से चढ़ता है (अधिक फ़ोनॉन)।
एक वाहक-गिनती का तापमापी है, दूसरा प्रकीर्णन का। (ग)
वाहक इलेक्ट्रॉन--विवर \emph{जोड़ों} में जन्मते हैं, और
संतुलन $np = n_{\text{i}}^2$ अंतराल की लागत दोनों में बाँट देता है
--- और इसीलिए $E_{\text{g}}/2$। (घ)
$n_{\text{i}}$ \qty{760}{K} के पास $\qty{5e21}{m^{-3}}$ तक पहुँच
जाता है: \omterm{def:b3:electrons-in-solids:doping}{अपमिश्रण} डूब जाता है और परिपथ अपनी अभिकल्पना भूल जाता है।
असली उपकरण उससे पहले ही हार मान लेते हैं --- उनके उत्क्रम रिसाव, जो
हर दस केल्विन पर दुगने होते हैं, पहले ही बिगड़ जाते हैं।
\end{solution}

\begin{solution}{exo:b3:electrons-in-solids:9}
(क) $\qty{5e22}{m^{-3}}$ --- यानी $n_{\text{i}}$ का पचास लाख गुना।
(ख) वही गुणक $\sim5\times10^6$: दस लाख परमाणुओं पर एक कण
पूरी चालकता का मालिक बन बैठता है। (ग) क्योंकि आकस्मिक भाग प्रति दस
लाख भी बिन बुलाए यही कर देते हैं: केवल भाग प्रति अरब तक शुद्ध किया
गया \omterm{def:b3:crystalline-solids:lattice}{क्रिस्टल} ही ऐसी चालकता देता है जो अभिकल्पक की अपनी हो। (घ)
$d \sim n^{-1/3}
\approx \qty{27}{nm}$, यानी \qty{2.4}{nm} की कक्षा का दस गुना: \omterm{def:b3:electrons-in-solids:doping}{दाता}
अलग-थलग हाइड्रोजन हैं; \omterm{def:b3:electrons-in-solids:doping}{अपमिश्रण} सौ गुना बढ़ाइए और वे
कक्षाएँ छूने लगती हैं --- तब कोई अशुद्धि-बैंड, और अंततः एक \omterm{def:b3:electrons-in-solids:classes}{धातु}।
\end{solution}

\begin{solution}{exo:b3:electrons-in-solids:10}
(क) $E = 13.6\times0.26/11.7^2 \approx \qty{26}{meV}$ (फ़ॉस्फ़ोरस का
मापा गया मान, \qty{45}{meV}, उस पैमाने को ईमानदार रखता है)। (ख)
$k_{\text{B}}T = \qty{26}{meV}$ के तुलनीय:
कमरे के ताप पर लगभग सारे \omterm{def:b3:electrons-in-solids:doping}{दाता} आयनित हैं --- और यही
\cref{def:b3:electrons-in-solids:doping} की मान्यता है। (ग) $a =
\qty{0.53}{\angstrom}\times\varepsilon_{\text{r}}/(m^*/
m_{\text{e}}) \approx \qty{2.4}{nm}$। (घ) दर्जनों कोष्ठिकाओं में
फैली कोई कक्षा उस \omterm{def:b3:crystalline-solids:lattice}{क्रिस्टल} को एक चिकने माध्यम की तरह देखती है ---
और ठीक यही वह शर्त है जिसके अंतर्गत कोई पिंडीय
$\varepsilon_{\text{r}}$ और कोई बैंड-औसत $m^*$ वैध हैं:
वह सन्निकटन स्वयं ही अपना प्रमाणपत्र है।
\end{solution}

\begin{solution}{exo:b3:electrons-in-solids:11}
(क) $V_{\text{bi}} = 0.0259\ln(10^{12}) \approx \qty{0.71}{V}$।
(ख) $eV/k_{\text{B}}T = 19.3$: अनुपात $\eu^{19.3} \approx
2\times10^{8}$ --- यानी वह डायोड आठ अंकों तक एकतरफ़ा रास्ता है।
(ग) $I_0 \propto n_{\text{i}}^2 \propto
\eu^{-E_{\text{g}}/k_{\text{B}}T}$: वही चरघातांकी, जो
\cref{exo:b3:electrons-in-solids:8} में था, और इसीलिए वह
अंगूठे का नियम कि वह दुगना हो जाता है। (घ) चारों डायोड एक हीरा बनाते
हैं: हर अर्ध-चक्र में जो जोड़ा अग्र-अभिनत हो वही धारा को भार से होकर
\emph{उसी} दिशा में मोड़ देता है --- भीतर AC, बाहर ऊबड़-खाबड़
DC।
\end{solution}

\begin{solution}{exo:b3:electrons-in-solids:12}
(क) अंतराल सँकरा कीजिए और अधिक फ़ोटॉन उसे पार कर लेते हैं, पर हर एक
केवल छोटी अंतराल-ऊर्जा देता है; चौड़ा कीजिए और हर फ़ोटॉन अधिक चुकाता
है पर कम ही योग्य ठहरते हैं: फ़सल $\times$ वोल्टता बीच में शिखर बनाती है।
(ख) \qty{1.3}{eV} के इष्टतम से ज़रा नीचे: सिलिकॉन की आदर्श
छत $\sim30\,\%$ है --- इतना पास कि उसका सस्तापन जीत जाता है।
(ग) चौड़े-अंतराल वाला ऊपरी सेल नीले को ऊँची वोल्टता पर लेता है और
लाल को नीचे सँकरे-अंतराल वाले सेल तक जाने देता है: हर फ़ोटॉन अपनी ही
ऊर्जा के पास बटोरा जाता है, तापीयकरण सिकुड़ जाता है, और उस चट्टे की
छत चालीसों की ओर चढ़ जाती है। (घ) ऊष्मा $I_0$ को
चरघातांकी ढंग से बढ़ाती है, जिससे खुला-परिपथ वोल्टता
$\sim(k_{\text{B}}T/e)\ln(I_{\text{L}}/I_0)$ नीचे खिंचती है; और
अंतराल स्वयं भी ज़रा सिकुड़ जाता है, जिससे वोल्टता के बदले जो धारा
मिलती है वह पूरी भरपाई नहीं कर पाती।
\end{solution}

\begin{solution}{pb:b3:electrons-in-solids:1}
\textbf{1.} $\lambda_{\text{शि}} \approx \qty{500}{nm}$:
यानी लगभग \qty{2.5}{eV}।
\textbf{2.} $hc/E_{\text{g}} \approx \qty{1.1}{\micro m}$।
\textbf{3.} वह सेलों के पार निकल जाता है (सोखने के लिए कोई अंतिम
अवस्था ही नहीं) और पिछली परत में ऊष्मा बनकर समाप्त होता है --- छत
गरम करता हुआ, तार नहीं।
\textbf{4.} $\qty{1.12}{eV}$ युग्म के रूप में बचती है; और शेष
\qty{1.4}{eV} पिकोसेकंडों में फ़ोनॉनों में बह जाती है --- किसी भी
परिपथ के रोक पाने से तेज़।
\textbf{5.} अंतराल-से-नीचे की पारदर्शिता (शक्ति का $\sim20\,\%$) और
अंतराल-से-ऊपर का तापीयकरण ($\sim30\,\%$): संधि के एक भी इलेक्ट्रॉन
देखने से पहले ही वर्णक्रम पर आधा कर लग जाता है।
\textbf{6.} अवशोषण एक एकल-फ़ोटॉन क्वांटम छलाँग है, जिसे कोई असली
अंतिम अवस्था चाहिए; आधे अंतराल पर ठहरने को कोई अवस्था है ही नहीं, और
सूर्य-प्रकाश के फ़ोटॉन-घनत्वों पर दो-फ़ोटॉन घटनाएँ नगण्य हैं।
\textbf{7.} \omterm{prop:b3:electrons-in-solids:junction}{अवक्षय क्षेत्र} में या उसके पास बने युग्म उस
अंतर्निर्मित क्षेत्र को अनुभव करते हैं: इलेक्ट्रॉन n ओर बुहार दिए
जाते हैं और \omterm{def:b3:electrons-in-solids:doping}{विवर} p ओर, और वह क्षेत्र उन्हें इतनी तेज़ी से अलग कर देता
है कि वे पुनर्योजन के लिए एक-दूसरे को ढूँढ़ ही नहीं पाते --- आवेश तब
तक जमा होता है जब तक कोई बाहरी परिपथ उसे धारा के रूप में राहत न दे।
\textbf{8.} $V_{\text{bi}} = 0.0259\ln(10^{12}) \approx
\qty{0.71}{V}$।
\textbf{9.} $60\times\qty{0.6}{V} = \qty{36}{V}$।
\textbf{10.} $I = 400/36 \approx \qty{11}{A}$।
\textbf{11.} पूरे पैनल पर अंतराल-से-ऊपर का प्रकाश
$\sim\qty{1400}{W}$ है, प्रति फ़ोटॉन $\sim\qty{2}{eV}$ पर: यानी प्रति
सेकंड $4\times10^{21}$ फ़ोटॉन --- पर वे साठ सेल \emph{श्रेणी} में
हैं, अतः वही \qty{11}{A} (यानी इलेक्ट्रॉन-अभिवाह $I/e \approx
7\times10^{19}\,\unit{s^{-1}}$) हर सेल से होकर गुज़रती है, और हर सेल
के हिस्से के फ़ोटॉन $4\times10^{21}/60 \approx
7\times10^{19}$ प्रति सेकंड हैं: यानी लगभग हर अवशोषित फ़ोटॉन एक
बटोरा गया इलेक्ट्रॉन देता है। क्वांटम दक्षता बढ़िया है; और
हानियाँ ऊर्जा की हैं, संख्या की नहीं।
\textbf{12.} वर्णक्रम के 50\,\% के बाद वह सेल वोल्टता की कमी
($0.6/1.12 \approx 54\,\%$) चुकाता है, और कुछ अंक
परावर्तन तथा पुनर्योजन के: $0.5\times0.54 \approx 27\,\%$,
जो छँटकर निर्धारित 20\,\% रह जाता है।
\textbf{13.} वही रिसाव $I_0$: खुला-परिपथ वोल्टता
$(k_{\text{B}}T/e)\ln(I_{\text{L}}/I_0)$ वहीं रुक जाती है जहाँ
प्रकाश-धारा और डायोड का रिसाव बराबर हों --- लगभग \qty{0.6}{V}, जो
अंतराल से काफ़ी कम है।
\textbf{14.} $\sin30^\circ = 0.5$: अकेली ज्यामिति से ही दोपहर की
रेटिंग का आधा --- बादलों से पहले ही।
\textbf{15.} $\Delta T = \qty{40}{K}$: $-16\,\%$, यानी
$\approx\qty{335}{W}$ बचते हैं --- सबसे धूप वाले दिन प्रति वॉट सबसे
अच्छे दिन नहीं होते।
\textbf{16.} ऊष्मा $I_0$ को चरघातांकी ढंग से फुला देती है, और
खुला-परिपथ वोल्टता --- किसी लघुगणक की झिड़की --- प्रति केल्विन
प्रतिशत के अंश जितनी गिर जाती है; और ज़रा सिकुड़ा हुआ अंतराल बाकी
काम पूरा कर देता है।
\textbf{17.} श्रेणी-तारबंदी एक ही उभयनिष्ठ धारा थोपती है: वह छाया
वाला सेल, जो कुछ नहीं बना रहा, अपने उनसठ साथियों से उत्क्रम अभिनति
में धकेल दिया जाता है और उनकी शक्ति को किसी तप्त धब्बे के रूप में
अपव्यय करता है --- वह लड़ी सबसे कमज़ोर सेल तक गला घोंट लेती है।
\textbf{18.} बाईपास डायोड धारा को उस छाया वाली उप-लड़ी के चारों ओर
एक अग्र पथ दे देता है, और पूरे पैनल के उत्पादन के बजाय केवल उसकी
वोल्टता की बलि चढ़ाता है।
\textbf{19.} प्रति वर्ष $\qty{400}{W}\times0.15\times\qty{8760}{h} \approx
\qty{530}{kWh}$।
\textbf{20.} प्रति वर्ष $\approx 130$ यूरो: प्रतिदाय मोटे
तौर पर $300/130 \approx 2.5$ वर्ष में, और उसके बाद दो दशकों का लाभ।
\textbf{21.} $500/530 \approx$ एक वर्ष: वह पैनल अपनी निर्माण-ऊर्जा
लगभग उतनी ही तेज़ी से लौटा देता है जितनी तेज़ी से अपनी क़ीमत।
\textbf{22.} कालापन कोई चुनाव नहीं, बल्कि एक \omterm{thm:b3:electrons-in-solids:bloch}{बैंड संरचना} है:
अवशोषण को किसी भरी अवस्था से एक फ़ोटॉन-ऊर्जा ऊपर कोई खाली अवस्था
चाहिए, और अंतराल से नीचे सिलिकॉन के पास देने को कोई है ही नहीं ---
रंग-रोगन अवस्थाएँ नहीं जोड़ सकता।
\textbf{23.} ऊपर कोई चौड़े-अंतराल वाला सेल चढ़ा दीजिए, जो नीले को
ऊँची वोल्टता पर ले और लाल को नीचे जाने दे: वह युगल एकल-संधि की छत
हरा देता है। पेच यह है: उस श्रेणी-चट्टे को परतों के बीच धाराएँ (और
क़ीमतें) मिलानी पड़ती हैं।
\textbf{24.} जो कुछ भी $\tau$ को छोटा करे और संधियों को ज़हर दे:
UV से बने दोष और भीतर विसरित अशुद्धियाँ वाहकों को प्रकीर्ण और
फँसा लेती हैं, नमी संपर्कों को खा जाती है, और तप्त धब्बे डायोडों को
बुढ़ा देते हैं --- यानी किसी ऐसे \omterm{def:b3:crystalline-solids:lattice}{क्रिस्टल} की धीमी एन्ट्रॉपी, जिससे
तीस वर्ष धूप में बैठने को कहा गया हो।
\textbf{25.} अंतराल सूर्य का आधा बटोरता है; संधि उन युग्मों को
\qty{36}{V} की क्रमित धारा में बदलती है; श्रेणी-लड़ियाँ,
छाया और गर्मी की तपिश अपना हिस्सा काटती हैं; और प्रति वर्ष
$\sim\qty{530}{kWh}$ के साथ, एक वर्ष के ऊर्जा-प्रतिदाय तथा
तीन वर्ष के धन-प्रतिदाय के सामने, मध्यस्थ का फ़ैसला है: छत ढँक दीजिए।
\end{solution}
